\section{The volume of the region with small least size}
\label{rv-section}

We prove the volume estimate needed in
\eqref{amb-rare-input}.  A sublattice of least normalized size
replaces the complete Harder--Narasimhan flag.  If that size is
less than one, the sublattice is totally isotropic.  Unfolding
the number of primitive isotropic sublattices of a fixed rank
then gives the required estimate.

The global classical inputs are the unitary Tamagawa-number
theorem, the Tamagawa number one of the special linear group,
the analytic class-number formula, the local-global principle
for Hermitian spaces, and Louboutin's residue bound.  We specify
their use and normalize all measures below.  We do not import
the submitted flag-volume estimate.

\subsection{Statement and the sublattice of least size}

Use the notation of Section~\ref{amb-section}.  Thus $E/F$ is
CM and unramified at every finite place,
$[F:\mathbb Q]=f$, $r_E=|D_E|^{1/(2f)}$, and $V$ has dimension
$n=p+q$ and signature $(p,q)$ at every real place.
The lattice $\Lambda_0$ is self-dual.  For
$x\in\mathcal X=G(F)\backslash G(\mathbb A_F)$ use the point
lattice and metrics in \eqref{amb-point-lattice}, and write
$\lambda_x$ for its least size \eqref{amb-size}.
Let $\mathbb P$ be Haar measure on $\mathcal X$, divided by
its total volume.

\begin{theorem}\label{rv-main}
Suppose
\[
       \log n\ge15,\qquad q^2\le n,\qquad
       \log r_E\le e^n .
\]
Then
\begin{equation}\label{rv-probability}
 \mathbb P\{x:\lambda_x<1/2\}
       \le \exp(-nf\log n/3).
\end{equation}
Consequently, for $b\ge2$, $g\ge1$, and
$\tau_0=(nfr_Egb)^{-11q}$, the complement of
$\{x:\lambda_x\ge\tau_0\}$ has probability at most
$\exp(-nf\log n/3)$.  This proves the weaker estimate
\eqref{amb-rare-input}.  No bound on $\log f$ is needed here.
\end{theorem}

All positive metrics below are compatible with the Hermitian
form: its matrix in an orthonormal basis has entries $1,-1$
on the diagonal.  A sublattice is called primitive if its
quotient in the ambient lattice is torsion-free.

\begin{lemma}\label{rv-minimum}
For every full $\mathcal O_E$-lattice $\Lambda$ the infimum
$\lambda(\Lambda)$ is attained by a nonzero primitive sublattice
$M$.  Put $s=s(M)$.  Every nonzero sublattice of $M$ has size
at least $s$, and every nonzero torsion-free quotient of $M$
has size at most $s$.
If $\Lambda$ is self-dual and $s<1$, then $M$ is totally
isotropic.  Its rank is therefore between one and $q$.
\end{lemma}
\begin{proof}
Saturation preserves rank and decreases covolume.  At a fixed
rank $j$, a primitive sublattice determines, up to sign, a
primitive decomposable vector in
$\bigwedge_{\mathbb Z}^{2fj}\Lambda$; its norm is the real
covolume.  A bounded ball contains only finitely many vectors
of this fixed exterior lattice.  Thus there are only finitely
many primitive sublattices of bounded covolume at each rank.
There are only $n$ possible ranks.  Taking, for example,
$s(\Lambda)$ as an initial upper bound proves attainment.

The assertion for sublattices follows from minimality.
For a primitive submodule $N\subset M$, of rank $a<j$,
orthogonal quotient metrics give
\[
                  H(M)=H(N)H(M/N).
\]
This identity follows by taking a basis of the underlying
$\mathbb Z$-lattice, extending a basis of $N$, and projecting
the remaining vectors orthogonally; the determinant becomes
block triangular.  The normalizing powers of $c_E$ multiply
because ranks add.  Since
$H(N)\ge s^{2fa}$ and $H(M)=s^{2fj}$, it gives
$s(M/N)\le s$.

Let $\overline{M}^{\,*}$ denote the conjugate
$\mathcal O_E$-linear dual, with its dual metrics.  The
Hermitian form restricted to $M$ defines a homomorphism
\[
                         T:M\longrightarrow\overline{M}^{\,*}.
\]
Its operator norm at each complex place is at most one:
it is the restriction of the unitary matrix $J_v$, followed
by orthogonal projection onto the dual of the subspace.
Suppose that $T\ne0$.  Let $Q=M/\ker T$ and $I=T(M)$, of
common positive rank $a$.  By the duality of sublattices and
quotients in Lemma~\ref{amb-size-duality}, the quotient
assertion just proved gives
\[
                 s(Q)\le s,\qquad s(I)\ge s^{-1}.
\]
The change-of-volume argument of
Lemma~\ref{amb-map-separation}, now with these domain and
target lattices, gives
\[
 s(I)\le s(Q)\Bigl(\prod_{v\mid\infty}
                                  \|T_v\|_{\rm op}\Bigr)^{1/f}.
\]
It would follow that $1\ge s^{-2}>1$, a contradiction.
Hence the restricted pairing is zero.  An isotropic subspace
of a Hermitian space of signature $(p,q)$ has dimension at
most $\min(p,q)$: projection to either definite summand is
injective on that subspace.  In particular its dimension
is at most $q$.
\end{proof}

For self-dual $\Lambda$, duality also gives $H(\Lambda)=1$.
If $q=0$, the lemma excludes $\lambda(\Lambda)<1$, while
$s(\Lambda)=1$ gives the reverse inequality.  Thus
$\lambda_x=1$ everywhere in that case.  In the rest of the
proof assume $q\ge1$.

\subsection{The quotient volumes and their Haar measures}

At a real place put on $\operatorname U(a,b)$ the Haar
measure whose density at the identity is Euclidean volume
for
\[
       \langle X,Y\rangle=\Re\operatorname{tr}(X^*Y)
\]
on its real Lie algebra.  We call this the
Hilbert--Schmidt normalization.  At finite places normalize
the stabilizer of a self-dual lattice to have volume one.
Let $\mathcal V_{a,b}$ be the resulting adelic quotient
volume, with signature $(a,b)$ at all real places, and set
$r=a+b>0$.
Write $d_F=|D_F|$ and
$L_\ell=L_F(\ell,\chi^\ell)$ for the factors in
Section~\ref{md-section}, where $\chi$ is the quadratic
character of $E/F$.

\begin{lemma}\label{rv-unitary-volume}
For $r>0$,
\begin{equation}\label{rv-unitary-volume-formula}
        \mathcal V_{a,b}
          =2d_F^{r^2/2}\prod_{\ell=1}^{r}L_\ell .
\end{equation}
Consequently, if $m=n-2j>0$, the ratio for signatures
$(p-j,q-j)$ and $(p,q)$ is
\begin{equation}\label{rv-volume-ratio}
 \frac{\mathcal V_{p-j,q-j}}{\mathcal V_{p,q}}
       =r_E^{-2fj(n-j)}
                      \prod_{\ell=m+1}^{n}L_\ell^{-1}.
\end{equation}
\end{lemma}
\begin{proof}
We use the classical unitary Tamagawa-number theorem in the
following form: over a fixed quadratic extension, the
Tamagawa quotient volume of a nondegenerate Hermitian
space is independent of its form and of its positive
dimension, and equals two in the commutative unitary case.
Weil proves the independence in \cite[\S54]{weil-siegel};
his unitary case has parameter $e=1/2$, as specified in
\cite[\S27]{weil-siegel}.  The threshold $r>4e-2$ is
therefore precisely $r>0$.  The absolute normalization
agrees with the unitary mass formula already stated in
\eqref{md-classical-mass}.

Here are the measure comparisons needed to apply this
theorem to our normalization.  At finite places the density
of the self-dual stabilizer is
$\prod_{\ell=1}^{r}(1-\eta_v^\ell q_v^{-\ell})$ by
\eqref{md-stabilizer-density}, independent of the Hermitian
form, including at dyadic places.  At infinity the conversion
from an algebraic invariant differential to
Hilbert--Schmidt volume is independent of the signature.
One can check the latter directly on the Lie algebra.
For a Hermitian matrix $H$, the map $X\mapsto HX$
identifies $X^*H+HX=0$ with the skew-Hermitian matrices.
After writing $H=AJA$ with $A$ positive definite and $J$
diagonal with entries $1,-1$, put $Y=AXA^{-1}$.
Then $HX=AJYA$.  Multiplication by $J$ is an isometry.
On the real space of skew-Hermitian matrices the map
$Z\mapsto AZA$ has determinant $|\det H|^r$:
diagonalize $A$, count its real diagonal coordinates and
its complex off-diagonal coordinates.  The corresponding
finite-place determinant is a unit in an integral
self-dual basis.  In arbitrary rational bases the
remaining determinant factors cancel by the product
formula.  Thus changing the signature introduces no
new constant in the normalized quotient volume.

It suffices to compute that volume for signature $(r,0)$.
In Hilbert--Schmidt normalization the compact group has
volume
\begin{equation}\label{rv-compact-volume}
 U_r:=\operatorname{vol}\operatorname U(r)
      =2^{r(r-1)/2}\prod_{\ell=1}^{r}
                        \frac{2\pi^\ell}{(\ell-1)!}.
\end{equation}
For completeness, in the fibration
$\operatorname U(r-1)\subset\operatorname U(r)\to S^{2r-1}$
the off-diagonal tangent directions have squared norm
twice their sphere norm; the remaining tangent direction
has unchanged norm.  The quotient metric volume is
$2^{r-1}$ times the standard sphere volume
$2\pi^r/(r-1)!$.  Induction, starting with the circle
of length $2\pi$, proves \eqref{rv-compact-volume}.
Multiplying the mass in Theorem~\ref{md-mass-ratio} by
$U_r^f$ cancels all factorials and powers of $2\pi$ and
gives \eqref{rv-unitary-volume-formula}.
Finally $d_F=r_E^f$ and
$n^2-(n-2j)^2=4j(n-j)$ prove \eqref{rv-volume-ratio}.
\end{proof}

We also need a general linear quotient volume.  At the
$f$ complex places use
\[
 da=\prod_{v\mid\infty}
              \frac{\prod_{\alpha,\beta}
                     d\Re A_{v,\alpha\beta}\,
                     d\Im A_{v,\alpha\beta}}
                   {|\det A_v|^{2j}},
 \qquad
 u(A)=\prod_{v\mid\infty}|\det A_v|^2.
\]
For a projective rank-$j$ $\mathcal O_E$-module $P$ embedded
as a lattice in $E^j$, let $\operatorname{Aut}(P)$ act on
$\operatorname{GL}_j(\mathbb C)^f$ by left multiplication,
and let $\mathcal G(P)$ be the volume of the part of that
quotient with $1\le u(A)<e$.
The number $u(A)$ is invariant under $\operatorname{Aut}(P)$,
since the absolute norm of the determinant of an integral
automorphism is one.

\begin{lemma}\label{rv-gl-volume}
With one representative of each isomorphism class of
projective rank-$j$ $\mathcal O_E$-modules,
\begin{equation}\label{rv-gl-formula}
 \mathcal G_j:=\sum_{[P]}\mathcal G(P)
     =c_E^{j^2}\operatorname{Res}_{s=1}\zeta_E(s)
                         \prod_{\ell=2}^{j}\zeta_E(\ell).
\end{equation}
For any nonnegative measurable function $h$,
\begin{equation}\label{rv-gl-radial}
 \int_{\operatorname{Aut}(P)\backslash
                         \operatorname{GL}_j(\mathbb C)^f}
                  h(u(A))\,da
      =\mathcal G(P)\int_0^\infty h(u)\,\frac{du}{u}.
\end{equation}
\end{lemma}
\begin{proof}
Here the classical input is the Tamagawa number one of
$\operatorname{SL}_j$ over a number field, recorded, for
example, in \cite[\S57]{weil-siegel}; for $j=1$ the special
linear group is trivial.  Together with the analytic
class-number formula for $E$, applied to the determinant
torus, it says that the norm-one adelic quotient of
$\operatorname{GL}_j$ has volume one with the following
normalization:
\begin{equation}\label{rv-tamagawa-gl}
 |D_E|^{-j^2/2}
       \bigl(\operatorname{Res}_{s=1}\zeta_E(s)\bigr)^{-1}
       \prod_{v\mid\infty}2^{j^2}\,da_v
       \prod_{w\nmid\infty}(1-q_w^{-1})^{-1}\,da_w.
\end{equation}
At a finite place, $da_w=|\det A|_w^{-j}dA$, with
$M_j(\mathcal O_{E,w})$ of additive volume one.
The measure on the norm-one group is obtained by
disintegration against $d\log|\det|_{\mathbb A_E}$.
To explain the compatibility, the invariant differential
on $\operatorname{GL}_j$ decomposes into the invariant
differential on its determinant-one fiber and $dt/t$.
The determinant sequence is split; applying Fubini to it
multiplies the special-linear Tamagawa number by the
idele-class norm-one volume.  The analytic class-number
formula gives the latter as one after the residue factor
in \eqref{rv-tamagawa-gl} is included.  The factor two in
each complex additive coordinate and the discriminant
factor are the usual additive adelic normalization.
This also checks \eqref{rv-tamagawa-gl} directly when $j=1$.

Counting ordered independent columns over the residue
field gives
\[
 \operatorname{vol}_{da_w}\operatorname{GL}_j(\mathcal O_{E,w})
                 =\prod_{\ell=1}^{j}(1-q_w^{-\ell}).
\]
The convergence factor in \eqref{rv-tamagawa-gl} cancels
the term with $\ell=1$.  Dividing out the finite compact
group and sorting the remaining components by the
associated projective module therefore gives
\[
 1=c_E^{-j^2}
       \bigl(\operatorname{Res}_{s=1}\zeta_E(s)\bigr)^{-1}
       \left(\prod_{\ell=2}^{j}\zeta_E(\ell)^{-1}\right)
       \sum_{[P]}\mathcal G(P).
\]
This proves \eqref{rv-gl-formula}; no assumption that
$P$ is free has been imposed.

For \eqref{rv-gl-radial}, multiplication at every complex
place by $u^{1/(2fj)}I_j$ gives a central one-parameter
subgroup with determinant norm $u$.  Haar disintegration
along this subgroup has radial measure $du/u$ and a
constant fiber volume.  The slab $1\le u<e$ has radial
volume one, so its volume $\mathcal G(P)$ is exactly that
constant.  This proves the formula, first for intervals
and then for all nonnegative measurable functions.
\end{proof}

\subsection{Primitive isotropic sublattices and the unipotent radical}

Fix a rank $j$ with $1\le j\le\min(p,q)$ and put
$m=n-2j>0$ and $d=2mj+j^2$.
Let $P$ be a primitive totally isotropic sublattice of a
self-dual Hermitian lattice $\Lambda$.

\begin{lemma}\label{rv-integral-splitting}
There is an integral orthogonal decomposition
\[
                     \Lambda=\mathbb H(P)\perp\Lambda',
\]
where $\mathbb H(P)$ is the hyperbolic module
$P\oplus\overline P^{\,*}$ and $\Lambda'$ is self-dual
of rank $m$ and signature $(p-j,q-j)$ at each real place.
The stabilizer of $P$ in $\operatorname{Aut}(\Lambda)$
is a semidirect product of its integral unipotent radical
with
$\operatorname{Aut}(P)\times\operatorname{Aut}(\Lambda')$.
\end{lemma}
\begin{proof}
Finite unramifiedness implies
$\operatorname{Tr}_{E/F}(\mathcal O_E)=\mathcal O_F$.
Indeed the local trace is surjective in an unramified
quadratic extension, and it is the sum map in the split
case.  The image of the global trace is an ideal of
$\mathcal O_F$ whose localization at every finite prime
is the unit ideal.  Choose $\theta\in\mathcal O_E$ with
$\theta+\bar\theta=1$.

Primitivity implies that $\Lambda/P$ is projective over
the Dedekind domain $\mathcal O_E$.  Restriction of the
perfect pairing on $\Lambda$ thus gives a surjection
$\Lambda\to\overline P^{\,*}$, and this surjection has
an $\mathcal O_E$-linear section.  The pairing of its
image with $P$ is perfect.  If $B$ is the Hermitian
Gram map of the section, subtract from the section
the map into $P$ corresponding to $\theta B$.
The new Gram map is
$B-\theta B-\bar\theta B=0$.
This calculation may be made after localization in
bases; it is intrinsic and hence holds for projective
modules.  We have obtained an integral isotropic
section, whose sum with $P$ is $\mathbb H(P)$.
Its pairing is perfect, so it is an orthogonal direct
summand.  Its orthogonal complement is the required
self-dual $\Lambda'$.

Every automorphism of $P$ extends to $\mathbb H(P)$
by its inverse conjugate dual, and every automorphism
of $\Lambda'$ extends by the identity.  Conversely,
an isometry preserving $P$ acts on $P$ and on
$P^\perp/P=\Lambda'$ by such automorphisms.
After removing these actions it lies in the unipotent
radical.  This proves the stated semidirect product.
The argument used a trace-one element, rather than
division by two, so it remains integral at dyadic places.
\end{proof}

At infinity choose hyperbolic coordinates
$E^j\oplus E^m\oplus E^j$ in which the Hermitian matrix is
\[
 \begin{pmatrix}0&0&I_j\\0&J'&0\\I_j&0&0\end{pmatrix},
 \qquad J'=\operatorname{diag}(1,\ldots,1,-1,\ldots,-1),
\]
and the compatible positive metric is the ordinary
Euclidean metric.  The unipotent radical consists of
\[
 n(Z,B)=
 \begin{pmatrix}
 I_j&-Z^*J'&B\\
 0&I_m&Z\\
 0&0&I_j
 \end{pmatrix},
                  \qquad B+B^*=-Z^*J'Z.
\]
Its tangent vectors are parametrized by $Z\in M_{m,j}(\mathbb C)$
and skew-Hermitian $C\in M_j(\mathbb C)$, and have
squared Hilbert--Schmidt norm
$2\|Z\|_{\rm HS}^2+\|C\|_{\rm HS}^2$.
Give the radical the Haar measure with this tangent density.

\begin{lemma}\label{rv-radical-covolume}
In the product of the radicals at the $f$ complex places,
the integral subgroup has covolume
\begin{equation}\label{rv-radical-formula}
                  r_E^{fd/2}H(P)^{n-j}.
\end{equation}
Here $H(P)$ is taken in the positive metrics just chosen.
\end{lemma}
\begin{proof}
Projection to $Z$ identifies the quotient by the center
with $\operatorname{Hom}_{\mathcal O_E}
(\overline P^{\,*},\Lambda')$.  It is surjective on
integral points: for any integral $Z$ one may take
$B=-\theta Z^*J'Z$.  The fibers are translates of the
integral skew-Hermitian central lattice.  Multiplication
in these coordinates is triangular with diagonal
Jacobian one, so its Haar measure is the product of
the two additive measures.  A fundamental domain may
be chosen by reducing $Z$ first and then the central
coordinate.  Its volume is the product of the two
additive covolumes, even though the group need not
be abelian.

For the $Z$ coordinate, the elementary determinant
identity for tensor products gives normalized height
\[
 H\bigl(\operatorname{Hom}(\overline P^{\,*},\Lambda')\bigr)
                =H(P)^m H(\Lambda')^j=H(P)^m.
\]
This identity follows in free bases from the determinant
of the Gram matrices and, for projective modules, by
localizing their inclusions in free lattices and
multiplying the resulting indices.
Its ordinary additive covolume is
$c_E^{mj}H(P)^m$.
The factor two in the squared $Z$-norm multiplies this
by $2^{fmj}$, giving $r_E^{fmj}H(P)^m$.

For clarity we also calculate the center, including its
diagonal coordinates.  Put
$\mathcal O_E^-=\{z\in\mathcal O_E:\bar z=-z\}$.
In the selected complex embeddings,
the real projection of $\mathcal O_E$ is
$\tfrac12\mathcal O_F$, by trace surjectivity.
Since $\operatorname{covol}(\mathcal O_E)
=2^{-f}d_F$, orthogonal projection gives
\[
              \operatorname{covol}(\mathcal O_E^-)
                         =\sqrt{d_F}=r_E^{f/2}.
\]
An off-diagonal skew-Hermitian entry has a second,
conjugate entry and hence has squared norm twice its
ordinary complex squared norm.  Its integral
covolume is $2^fc_E=r_E^f$.  For $P=\mathcal O_E^j$
the total central covolume is therefore
$r_E^{f j^2/2}$.

For general $P$, write it, in an $E$-basis, as
$\bigoplus_i\mathfrak a_i$.  An off-diagonal coordinate
belongs to $\mathfrak a_i\bar{\mathfrak a}_k$ and
contributes the additional factor
$N\mathfrak a_i\,N\mathfrak a_k$.
The invariant ideal $\mathfrak a_i\bar{\mathfrak a}_i$
is $\mathfrak b_i\mathcal O_E$ with
$N_{F/\mathbb Q}\mathfrak b_i=N_{E/\mathbb Q}\mathfrak a_i$.
This follows prime by prime, both in the split and in
the unramified inert case.  Its diagonal skew part is
$\mathfrak b_i\mathcal O_E^-$ and contributes
$N\mathfrak a_i$.  Thus the total ideal factor is
$(\prod_iN\mathfrak a_i)^j$.
For an arbitrary positive metric on $P$, a complex
change of basis $A$ acts on the center by $C\mapsto ACA^*$,
whose real determinant is $|\det A|^{2j}$.
Multiplication over the complex places supplies exactly
the remaining metric factor in $H(P)^j$.
The central covolume is
$r_E^{fj^2/2}H(P)^j$.
Multiplying it by the $Z$ covolume proves
\eqref{rv-radical-formula}, since $m+j=n-j$.
\end{proof}

\subsection{The real Iwasawa factor}

Put
\[
 S_a=\prod_{\ell=1}^{a}\frac{2\pi^\ell}{(\ell-1)!},
                   \qquad S_0=1.
\]
The stabilizer of the isotropic $j$-plane has Levi
subgroup
\[
       \operatorname{GL}_j(\mathbb C)\times
                          \operatorname U(p-j,q-j).
\]
Its embedding is
$(A,b)\mapsto\operatorname{diag}(A,b,A^{-*})$.
Let $K=\operatorname U(p)\times\operatorname U(q)$
carry probability Haar measure.

\begin{lemma}\label{rv-iwasawa}
For the Hilbert--Schmidt measures on the unitary group,
its radical, and the smaller unitary group, and for the
ordinary complex matrix Haar measure $da$ used above,
\begin{equation}\label{rv-iwasawa-formula}
 \begin{split}
 \int_{\operatorname U(p,q)}F(g)\,dg
   =c_j\,2^{j^2}
       \int_N\int_{\operatorname{GL}_j(\mathbb C)}
       \int_{\operatorname U(p-j,q-j)}\int_K
       &F\bigl(n\operatorname{diag}(A,b,A^{-*})k\bigr)\\[-1mm]
       &\hspace{-18mm}\cdot|\det A|^{-2(n-j)}
                         \,dk\,db\,da\,dn ,
 \end{split}
\end{equation}
for every nonnegative measurable $F$, where
\begin{equation}\label{rv-iwasawa-constant}
 c_j=
 2^{-d/2}\frac{U_pU_q}
           {2^{j^2/2}U_jU_{p-j}U_{q-j}}
 =2^{-j(j+1)/2}
                \frac{S_pS_q}{S_jS_{p-j}S_{q-j}} .
\end{equation}
\end{lemma}
\begin{proof}
The compact group $K$ is transitive on isotropic
$j$-planes: the projections of an orthonormal isotropic
frame to its positive and negative summands are scaled
orthonormal frames with the same scale; extend each
to a unitary basis.  This gives $G=NMK$.
The redundancy is the compact subgroup
$K_M\simeq\operatorname U(j)\times
\operatorname U(p-j)\times\operatorname U(q-j)$.
On its first factor the embedding duplicates a tangent
matrix, so its metric is twice the ordinary
Hilbert--Schmidt metric.  Its volume is
$2^{j^2/2}U_jU_{p-j}U_{q-j}$.
Likewise the full embedded general linear factor has
twice its usual tangent metric, and real dimension
$2j^2$, accounting for $2^{j^2}$ in
\eqref{rv-iwasawa-formula}.

It remains to compute the transverse Jacobian.
The radical Lie algebra is orthogonal to its Cartan
involutive image.  Its projection onto the orthogonal
complement of the compact Lie algebra is
$(X-\vartheta X)/2$, of norm $\|X\|/\sqrt2$.
It is orthogonal to the noncompact Levi directions.
There are $d$ radical directions, so this projection
contributes $2^{-d/2}$.
The quotient of the compact volumes is $U_pU_q/\operatorname{vol}K_M$,
giving the first expression for $c_j$.

Finally conjugation by the Levi acts on the complex
$Z$ coordinate with real determinant
$|\det A|^{2m}$ and on the skew-Hermitian central
coordinate with determinant $|\det A|^{2j}$.
Its modular character is consequently
$|\det A|^{2(n-j)}$.
Left invariance therefore requires its reciprocal in
the $NMK$ integral.  This determines the measure
globally from the just computed density at the identity.
Inserting \eqref{rv-compact-volume} into the first
expression in \eqref{rv-iwasawa-constant} gives the
second, by cancellation of the powers of two.
\end{proof}

\subsection{Unfolding the number of isotropic sublattices}

Let $N_j(x,T)$ count the primitive totally isotropic
rank-$j$ sublattices of $\Lambda_x$ whose size in its
point metrics is less than $T$.  This number is finite:
the exterior-lattice argument in Lemma~\ref{rv-minimum}
gives finiteness for any bound on the covolume.
It is measurable, being a countable sum of indicators
of continuous height inequalities on the arithmetic
components of the adelic quotient.

\begin{proposition}\label{rv-mean-value}
For $1\le j\le\min(p,q)$ and $m=n-2j>0$,
\begin{equation}\label{rv-mean-value-formula}
 \int_{\mathcal X}N_j(x,T)\,d\mathbb P(x)
       =\frac{\kappa_j}{n-j}\,
                        T^{2fj(n-j)}\qquad(T>0),
\end{equation}
where
\begin{equation}\label{rv-kappa}
 \kappa_j=
 r_E^{-fd/2}\,2^{-fj(j+1)/2}
 \left(\frac{S_pS_q}{S_{p-j}S_{q-j}S_j}\right)^f
 \operatorname{Res}_{s=1}\zeta_E(s)
 \frac{\prod_{\ell=2}^{j}\zeta_E(\ell)}
      {\prod_{\ell=m+1}^{n}L_\ell},
 \qquad d=2mj+j^2 .
\end{equation}
\end{proposition}
\begin{proof}
We give the orbit and measure details of the unfolding.
After division by $K_f$, the adelic quotient is the
disjoint union of
$\operatorname{Aut}(\Lambda)\backslash G_\infty$,
over the isometry classes in the genus of $\Lambda_0$.
Finite compact volume is one, so the measures on these
components are precisely the archimedean quotient
measures.  For each component, split the sum defining
$N_j$ into orbits under $\operatorname{Aut}(\Lambda)$.
The usual change of variables in a fundamental domain
gives
\[
 \int_{\operatorname{Aut}(\Lambda)\backslash G_\infty}
                    \sum_P {\bf1}_{s_x(P)<T}\,dx
  =\sum_{[P]}
       \int_{\operatorname{Aut}(\Lambda,P)\backslash G_\infty}
                                  {\bf1}_{s_x(P)<T}\,dx .
\]
Here the sum on the right is over primitive isotropic
sublattices modulo the indicated automorphism group.
All summands are nonnegative, so Tonelli justifies
this equality before integrability is known.

The isometry classes of marked pairs $(\Lambda,P)$
are exactly the pairs consisting of an isomorphism
class of projective rank-$j$ modules $P$ and an
isometry class of self-dual lattices $\Lambda'$ of
signature $(p-j,q-j)$.
Indeed Lemma~\ref{rv-integral-splitting} gives this
description and its converse by
$\mathbb H(P)\perp\Lambda'$.
For the converse to lie in the prescribed rational
space, we use the classical local-global principle
for Hermitian spaces
\cite[Theorem 2.4.1]{kirschmer-herm}.
At every finite place all unimodular Hermitian
lattices over the unramified quadratic algebra are
isometric: first use the finite-field classification
in Lemma~\ref{md-unitary-order}, and then the successive
Gram-matrix lifting in Proposition~\ref{md-lifting}.
The split case is the corresponding perfect-pairing
calculation.  At infinity the signatures agree.
Thus the rational spaces are isometric and the
resulting lattice lies in the same genus.
An isometry of marked pairs induces an isomorphism
of $P$ and an isometry on $P^\perp/P$; conversely
these extend to the hyperbolic orthogonal sums.
No class is lost by requiring a free basis for $P$,
because no such requirement has been made.

For one marked pair, its stabilizer is the semidirect
product in Lemma~\ref{rv-integral-splitting}.  Apply
\eqref{rv-iwasawa-formula} at the $f$ complex places.
Its radical quotient contributes
$r_E^{fd/2}H(P)^{n-j}$ by
Lemma~\ref{rv-radical-covolume}.  Denote by
$V(\Lambda')$ the volume of
$\operatorname{Aut}(\Lambda')\backslash
\operatorname U(p-j,q-j)^f$.
If $u=\prod_v|\det A_v|^2$, the size of the fixed
$P$ in the point metric is
\[
                         s_x(P)=(H(P)/u)^{1/(2fj)}.
\]
Indeed $n$ acts trivially on $P$, the Levi acts there
as $A$, and the right compact factor preserves
the positive norm.  Thus the integral of this
marked pair equals
\[
 \begin{split}
 c_j^f2^{fj^2}\,r_E^{fd/2}H(P)^{n-j}V(\Lambda')
       \int_{\operatorname{Aut}(P)\backslash
                           \operatorname{GL}_j(\mathbb C)^f}
       u^{-(n-j)}
       {\bf1}_{u>H(P)T^{-2fj}}\,da .
 \end{split}
\]
This also describes a fundamental domain in the
semidirect quotient: first reduce the Levi
coordinate and then the radical coordinate.
Integral Levi automorphisms preserve the radical
lattice and have global modular character one.
Consequently the radical covolume just used has no
extra index or stabilizer factor.

By \eqref{rv-gl-radial}, the last integral is
\[
 \mathcal G(P)\int_{H(P)T^{-2fj}}^\infty
                                u^{-(n-j)}\,\frac{du}{u}
  =\frac{\mathcal G(P)}{n-j}
                    H(P)^{-(n-j)}T^{2fj(n-j)}.
\]
The power of $H(P)$ cancels exactly.  Summing over
the two sets of lattice classes and dividing by
$\mathcal V_{p,q}$ proves
\[
 \int N_j(x,T)\,d\mathbb P(x)
 =c_j^f2^{fj^2}r_E^{fd/2}\mathcal G_j
      \frac{\mathcal V_{p-j,q-j}}{\mathcal V_{p,q}}
                  \frac{T^{2fj(n-j)}}{n-j}.
\]
Inserting \eqref{rv-gl-formula},
\eqref{rv-volume-ratio}, and
\eqref{rv-iwasawa-constant}, and using
$c_E=(r_E/2)^f$, yields \eqref{rv-kappa}.
In particular the computed integral is finite,
justifying all ordinary integral assertions as well.
\end{proof}

\subsection{An elementary bound for the constant}

\begin{lemma}\label{rv-kappa-bound}
Under the hypotheses of Theorem~\ref{rv-main},
for every $1\le j\le q$,
\begin{equation}\label{rv-kappa-small}
                 \kappa_j\le
                    \exp(-jnf\log n/2).
\end{equation}
\end{lemma}
\begin{proof}
The inequalities imply $p=n-q\ge q$ and
$n-2j\ge n-2\sqrt n\ge2$, so all preceding formulas
apply.  The first two powers in \eqref{rv-kappa}
are at most one, since $r_E\ge1$.

The factorials give a particularly large saving.
Direct cancellation shows
\[
 \frac{S_q}{S_{q-j}S_j}
   =\pi^{j(q-j)}
          \prod_{a=1}^{j}\frac{(a-1)!}{(q-j+a-1)!}
   \le\pi^{j(q-j)}.
\]
Moreover
$S_p/S_{p-j}\le[2\pi^p/(p-j)!]^j$.
Put $h=p-j\ge n-2\sqrt n$.  Hence
\[
 \log\frac{S_pS_q}{S_{p-j}S_{q-j}S_j}
       \le j\bigl(\log2+(n-j)\log\pi-\log(h!)\bigr).
\]
The integral of $\log x$ from one to $h$ gives
$\log(h!)\ge h\log h-h+1$.
Since $\log2<1$, $\log\pi<2$, and $h\le n$,
the last displayed right side is at most
$j(3n-h\log h)$.
For $\sqrt n\ge64$,
\[
 \log h\ge\log n+\log(1-2/\sqrt n)
                          \ge\log n-4/\sqrt n .
\]
Thus
\[
 3n-h\log h
 \le -n\log n+2\sqrt n\log n+4\sqrt n+3n .
\]
The ratio of the last three terms to $n\log n$
is at most
$1/32+1/240+1/5<1/4$, using $\log n\ge15$.
Consequently
\begin{equation}\label{rv-sphere-saving}
 \log\frac{S_pS_q}{S_{p-j}S_{q-j}S_j}
                         \le-\frac34jn\log n.
\end{equation}

For the residue we use the following unconditional
classical bound, valid for every number field $K$
of degree $d>1$:
\[
 \operatorname{Res}_{s=1}\zeta_K(s)
       \le\left(\frac{e\log|D_K|}{2(d-1)}\right)^{d-1}.
 \qquad\text{\cite[Theorem 1(1)]{louboutin-residue}}
\]
Here $d=2f$ and $\log|D_E|=2f\log r_E$, so it gives
\begin{equation}\label{rv-residue-bound}
       \log\operatorname{Res}_{s=1}\zeta_E(s)
                           \le 2f(n+1).
\end{equation}
Indeed $f/(2f-1)\le1$ and
$\log\log r_E\le n$.  The discriminant is greater
than one for a CM field, so these logarithms are
defined; if the logarithm of the residue bound is
negative, the same upper estimate still holds.

For completeness, the remaining Euler products
require only
\[
 \sum_{\ell=2}^{\infty}\log\zeta(\ell)
 \le\sum_{\ell=2}^{\infty}(\zeta(\ell)-1)
 =\sum_{a=2}^{\infty}\frac1{a(a-1)}=1.
\]
Since $\zeta_E(\ell)\le\zeta(\ell)^{2f}$, the numerator
in \eqref{rv-kappa} contributes at most $2f$ to its
logarithm.  Each inverse factor $L_\ell^{-1}$ is at
most $\zeta_F(\ell)\le\zeta(\ell)^f$: compare
$1-\eta_v^\ell q_v^{-\ell}$ with
$(1-q_v^{-\ell})^{-1}$ at every finite prime.
The denominator therefore contributes at most $f$.
Combining these estimates with
\eqref{rv-sphere-saving} and \eqref{rv-residue-bound},
\[
        \log\kappa_j
          \le-\frac34jnf\log n+2fn+5f
          \le-\frac12jnf\log n .
\]
The last inequality follows from
$\tfrac14jn\log n\ge\tfrac{15}{4}n\ge2n+5$.
\end{proof}

\begin{proof}[Proof of Theorem~\ref{rv-main}]
The case $q=0$ was already proved.
For $q\ge1$, if $\lambda_x<1/2$,
Lemma~\ref{rv-minimum} gives a primitive isotropic
sublattice of some rank $1\le j\le q$ and size
less than $1/2$.  Hence, pointwise,
\[
              {\bf1}_{\lambda_x<1/2}
                     \le\sum_{j=1}^{q}N_j(x,1/2).
\]
Put $A=nf\log n$.  The mean-value formula and
\eqref{rv-kappa-small} give
\[
 \mathbb P\{\lambda_x<1/2\}
 \le\sum_{j=1}^{q}
        \frac{2^{-2fj(n-j)}}{n-j}\kappa_j
 \le\sum_{j\ge1}e^{-jA/2}
 \le2e^{-A/2}\le e^{-A/3}.
\]
Here $A\ge15$, so both final inequalities hold.
This proves \eqref{rv-probability}.
If $q\ge1$, then $nfr_Egb\ge2n>2$ and
$\tau_0<1/2$; its exceptional set is smaller.
If $q=0$, then $\tau_0=1$ and $\lambda_x=1$,
so the exceptional set is empty.
\end{proof}

\paragraph{Flags and remaining estimates.}
For a flag with one nonzero rank increment, the constant
\eqref{rv-kappa} agrees with the normalization
in~\cite{alpo-uniform-source}.  The existence of a lattice
whose largest quotient size is less than $1/2$ is equivalent
to $\lambda_x<1/2$.  In one direction, such a lattice itself
has size less than $1/2$.  In the other, the minimizing lattice
in Lemma~\ref{rv-minimum} has all quotient sizes at most its size.
Thus \eqref{rv-probability} also bounds the probability
that such a nonzero lattice exists.
This section settles the rare-volume input only.
The integrated count of nearby nonidentity unipotents,
the deep-cusp mass bound, and the finite-rank kernel
identity in the ambient comparison remain separate.


% A direct proof of the integrated count of nearby rational isometries.
% Labels in this fragment have the prefix nu-.
